\documentclass[12pt]{article}
\usepackage[utf8]{inputenc}
\usepackage[T1]{fontenc}
\usepackage{hyperref}
\title{Wamas Language Resources for K-12 Curriculum Development}
\author{Compiled by Claude Sonnet 5.5 (claude-sonnet-5-5) for Andr\'as Kornai}
\date{}
\begin{document}
\maketitle

\section{Introduction}
Wamas is a Madang language of the Trans-New Guinea family, spoken on the northern coast of Madang Province, Papua New Guinea, about 40 km north of Madang town. The Endangered Languages Project lists about 220 speakers (2000) in 10 villages. Very little documentation turned up in searches.

\section{Primary Educational Resources}
I found no primers, textbooks or school programmes for Wamas in search results.

\section{Linguistic and Reference Materials}
\begin{itemize}
\item Glottolog classification (Wamas, wama1247): \url{https://glottolog.org/resource/languoid/id/wama1247.bigmap.html}
\item Wikipedia overview: \url{https://en.wikipedia.org/wiki/Wamas_language}
\end{itemize}
I saw no grammar, dictionary or wordlist for Wamas.

\section{Religious and Community Resources}
Joshua Project reports that Bible translation work in Wamas has begun, but I saw no published scripture: \url{https://joshuaproject.net/index.php/people_groups/15844}

\section{Digital and Audio Resources}
The Endangered Languages Project has an entry for Wamas, but I could not see the contents of its resources section: \url{https://endangeredlanguages.com/lang/3335}

\section{Curriculum Development Context}
There is almost no published material, so a curriculum would need new documentation work with the community. This should be done in partnership with the speakers of the ten villages.

\end{document}